\heading{理论物理基础（II）-模拟题4-期中}

\noteinfo{题目和答案由AI生成。}

\textbf{计算题（共 100 分）}

\begin{question}[16分]
有 $N$ 个彼此可区分、互不相互作用的粒子，每个粒子只能取三个单粒子能级
\[
\varepsilon_0=0,\qquad \varepsilon_1=\varepsilon,\qquad \varepsilon_2=2\varepsilon.
\]
设占据数分别为 $n_0,n_1,n_2$，总能量为 $U$。
\begin{enumerate}[label=（\arabic*）,leftmargin=2.8em,itemsep=0.2em]
\item 写出给定 $(n_0,n_1,n_2)$ 时的多重度 $\Omega$ 及约束条件；
\item 在 $N\gg1$ 时，用最大多重度原理求最概然占据数 $n_r^\ast/N$；
\item 特别地，若 $U=N\varepsilon$，求最概然的 $(n_0,n_1,n_2)$。
\end{enumerate}
\begin{solution}

由于粒子可区分，给定占据数时的多重度为
\[
\Omega=\frac{N!}{n_0!\,n_1!\,n_2!}.
\]
约束条件为
\[
n_0+n_1+n_2=N,\qquad
n_1+2n_2=\frac{U}{\varepsilon}.
\]

在大数极限下，令
\[
S=k\ln\Omega
\approx k\left(N\ln N-\sum_{r=0}^2 n_r\ln n_r\right).
\]
对
\[
\Phi=\ln\Omega-\alpha(n_0+n_1+n_2-N)-\beta(\varepsilon n_1+2\varepsilon n_2-U)
\]
分别对 $n_0,n_1,n_2$ 求偏导并令零，可得
\[
\ln n_r = \text{常数}-\beta \varepsilon_r.
\]
因此
\[
n_r^\ast = A e^{-\beta\varepsilon_r},
\]
由归一化条件得到
\[
\frac{n_r^\ast}{N}=
\frac{e^{-\beta\varepsilon_r}}{1+e^{-\beta\varepsilon}+e^{-2\beta\varepsilon}}.
\]
即
\ansmath{
\frac{n_0^\ast}{N}=\frac{1}{1+e^{-\beta\varepsilon}+e^{-2\beta\varepsilon}},\quad
\frac{n_1^\ast}{N}=\frac{e^{-\beta\varepsilon}}{1+e^{-\beta\varepsilon}+e^{-2\beta\varepsilon}},\quad
\frac{n_2^\ast}{N}=\frac{e^{-2\beta\varepsilon}}{1+e^{-\beta\varepsilon}+e^{-2\beta\varepsilon}}.
}

若 $U=N\varepsilon$，则平均每粒子能量为 $\varepsilon$，即
\[
\frac{n_1^\ast+2n_2^\ast}{N}=1.
\]
设 $x=e^{-\beta\varepsilon}$，则
\[
\frac{x+2x^2}{1+x+x^2}=1
\Longrightarrow x^2=1.
\]
取正根 $x=1$，所以 $\beta=0$，即最概然分布为三能级等占据：
\ansmath{
n_0^\ast=n_1^\ast=n_2^\ast=\frac{N}{3}.
}
\end{solution}
\end{question}

\begin{question}[16分]
四个粒子可占据如下单粒子能级：一个能量为 $0$ 的态、两个简并的能量为 $\varepsilon$ 的态、两个简并的能量为 $2\varepsilon$ 的态。求总能量恰为 $2\varepsilon$ 时，多体微观态数在以下三种情形下各是多少：
\begin{enumerate}[label=（\arabic*）,leftmargin=2.8em,itemsep=0.2em]
\item 四个粒子彼此可区分（Maxwell--Boltzmann 计数）；
\item 四个无相互作用玻色子；
\item 四个无自旋费米子。
\end{enumerate}
\begin{solution}

总能量为 $2\varepsilon$ 时，可能的能量分配只有两类：

\begin{enumerate}[label=（\arabic*）,leftmargin=2.4em,itemsep=0.2em]
\item 两个粒子处于能量 $\varepsilon$ 的态，其余两个处于能量 $0$ 的态；
\item 一个粒子处于能量 $2\varepsilon$ 的态，其余三个处于能量 $0$ 的态。
\end{enumerate}

对 Maxwell--Boltzmann 计数：

第一类中，先选出哪两个可区分粒子具有能量 $\varepsilon$，有
\[
\binom{4}{2}=6
\]
种；再给这两个粒子分别指定两个简并 $\varepsilon$ 态中的哪一个，共有 $2^2=4$ 种，所以第一类共有
\[
6\times4=24
\]
种。

第二类中，先选出哪个粒子在 $2\varepsilon$ 能级，有 $4$ 种；再选其所在的简并态，有 $2$ 种，所以第二类共有
\[
4\times2=8
\]
种。

故 Maxwell--Boltzmann 总微观态数为
\ansmath{\Omega_{\mathrm{MB}}=24+8=32.}

对玻色子计数：

第一类中，能量为 $\varepsilon$ 的两个玻色子分配到两个简并 $\varepsilon$ 态的方法数为
\[
\binom{2+2-1}{2}=\binom32=3.
\]
第二类中，唯一的一个 $2\varepsilon$ 玻色子可以处于两个简并态中的任一个，因此有 $2$ 种。故
\ansmath{\Omega_{\mathrm{B}}=3+2=5.}

对无自旋费米子计数：

系统共有 $1+2+2=5$ 个单粒子态，而有 $4$ 个费米子。由于泡利不相容原理，每态至多一个粒子。此时四个费米子的最低可能总能量已经是
\[
0+\varepsilon+\varepsilon+2\varepsilon=4\varepsilon.
\]
因此总能量 $2\varepsilon$ 根本不可能实现，所以
\ansmath{\Omega_{\mathrm{F}}=0.}
\end{solution}
\end{question}

\begin{question}[16分]
设一个 Einstein 晶体由 $N$ 个相互独立、频率都为 $\omega$ 的量子谐振子组成。忽略零点能。
\begin{enumerate}[label=（\arabic*）,leftmargin=2.8em,itemsep=0.2em]
\item 求单振子配分函数 $z$ 及全系统配分函数 $Z_N$；
\item 求平均能量 $U$；
\item 求定体热容 $C_V$，并给出高温与低温极限。
\end{enumerate}
\begin{solution}

单振子的能级为
\[
0,\ \hbar\omega,\ 2\hbar\omega,\ \cdots
\]
故单振子配分函数是几何级数
\[
z=\sum_{n=0}^{\infty}e^{-n\beta\hbar\omega}
=\frac{1}{1-e^{-\beta\hbar\omega}}.
\]
由于各振子独立，
\ansmath{
Z_N=z^N=\left(1-e^{-\beta\hbar\omega}\right)^{-N}.
}

平均能量
\[
U=-\frac{\partial}{\partial\beta}\ln Z_N
=-N\frac{\partial}{\partial\beta}\ln\left(1-e^{-\beta\hbar\omega}\right)^{-1}
= N\frac{\hbar\omega}{e^{\beta\hbar\omega}-1}.
\]
因此
\ansmath{
U=N\frac{\hbar\omega}{e^{\hbar\omega/(kT)}-1}.
}

令
\[
x=\frac{\hbar\omega}{kT},
\]
则
\[
C_V=\left(\frac{\partial U}{\partial T}\right)_V
=Nk\,\frac{x^2e^x}{(e^x-1)^2}.
\]
故
\ansmath{
C_V=Nk\frac{\left(\hbar\omega/kT\right)^2e^{\hbar\omega/(kT)}}{\left(e^{\hbar\omega/(kT)}-1\right)^2}.
}

高温极限 $x\ll1$ 时，
\[
e^x-1\approx x,
\]
故
\ansmath{C_V\approx Nk.}

低温极限 $x\gg1$ 时，
\[
e^x-1\approx e^x,
\]
故
\ansmath{
C_V\approx Nk\,x^2e^{-x}
=Nk\left(\frac{\hbar\omega}{kT}\right)^2e^{-\hbar\omega/(kT)}.
}
\end{solution}
\end{question}

\begin{question}[16分]
体积为 $V=V_1+V_2$ 的容器中有 $N$ 个彼此可区分但物理上完全相同的经典理想气体分子，体系处于平衡态。用组合计数方法求：
\begin{enumerate}[label=（\arabic*）,leftmargin=2.8em,itemsep=0.2em]
\item 左侧子体积 $V_1$ 中恰有 $n$ 个粒子的概率 $P(n)$；
\item $\langle n\rangle$ 与 $\mathrm{Var}(n)$；
\item 当 $V_1=V_2$ 时，最概然值与相对涨落的量级。
\end{enumerate}
\begin{solution}

平衡时每个粒子独立地以概率
\[
p=\frac{V_1}{V}
\]
出现在左侧，以概率 $1-p$ 出现在右侧。因此左侧恰有 $n$ 个粒子的概率服从二项分布
\[
P(n)=\binom{N}{n}p^n(1-p)^{N-n}.
\]
于是
\ansmath{
P(n)=\binom{N}{n}\left(\frac{V_1}{V}\right)^n\left(\frac{V_2}{V}\right)^{N-n}.
}

由二项分布性质，
\[
\langle n\rangle=Np=N\frac{V_1}{V},
\qquad
\mathrm{Var}(n)=Np(1-p)=N\frac{V_1V_2}{V^2}.
\]
故
\ansmath{
\langle n\rangle=N\frac{V_1}{V},
\qquad
\mathrm{Var}(n)=N\frac{V_1V_2}{V^2}.
}

当 $V_1=V_2=V/2$ 时，
\[
p=\frac12,
\qquad
n_{\mathrm{mp}}\approx \frac{N}{2},
\qquad
\mathrm{Var}(n)=\frac{N}{4}.
\]
于是标准差
\[
\sigma_n=\frac{\sqrt N}{2},
\]
相对涨落为
\[
\frac{\sigma_n}{\langle n\rangle}=\frac{1}{\sqrt N}.
\]
因此
\ansmath{
n_{\mathrm{mp}}\approx \frac{N}{2},
\qquad
\frac{\sigma_n}{\langle n\rangle}\sim N^{-1/2}.
}
\end{solution}
\end{question}

\begin{question}[16分]
经典理想气体的巨配分函数可以写成
\[
\Xi=\sum_{N=0}^{\infty} z^N Z_N,
\qquad
Z_N=\frac{1}{N!}\left(\frac{V}{\lambda^3}\right)^N,
\]
其中 $z=e^{\beta\mu}$，$\lambda$ 为热德布罗意波长。
\begin{enumerate}[label=（\arabic*）,leftmargin=2.8em,itemsep=0.2em]
\item 求 $\Xi$；
\item 求粒子数分布 $P(N)$；
\item 求 $\langle N\rangle$、$\mathrm{Var}(N)$，并说明这一结果与上一题的二项分布在稀薄极限下如何联系。
\end{enumerate}
\begin{solution}

由题给
\[
\Xi=\sum_{N=0}^{\infty}\frac{1}{N!}\left(z\frac{V}{\lambda^3}\right)^N
=\exp\!\left(z\frac{V}{\lambda^3}\right).
\]
故
\ansmath{
\Xi=\exp\!\left(z\frac{V}{\lambda^3}\right).
}

粒子数分布为
\[
P(N)=\frac{z^N Z_N}{\Xi}
=\frac{1}{N!}\left(z\frac{V}{\lambda^3}\right)^N
\exp\!\left(-z\frac{V}{\lambda^3}\right).
\]
设
\[
\bar N=z\frac{V}{\lambda^3},
\]
则
\ansmath{
P(N)=e^{-\bar N}\frac{\bar N^N}{N!},
}
即泊松分布。

泊松分布满足
\[
\langle N\rangle=\bar N,\qquad
\mathrm{Var}(N)=\bar N.
\]
因此
\ansmath{
\langle N\rangle=\mathrm{Var}(N)=z\frac{V}{\lambda^3}.
}

上一题的二项分布在极限
\[
N\to\infty,\qquad p\to0,\qquad Np=\text{常数}
\]
下趋于泊松分布，因此巨正则系综中“小子系统的粒子数涨落为泊松分布”正可看作固定总粒子数系统在稀薄极限下的局域结果。
\end{solution}
\end{question}

\begin{question}[20分]
考虑二维面积为 $A$ 的理想费米气体，粒子自旋为 $1/2$，质量为 $m$，温度取 $T=0$。
\begin{enumerate}[label=（\arabic*）,leftmargin=2.8em,itemsep=0.2em]
\item 求二维态密度 $g(\varepsilon)$；
\item 求费米能 $\varepsilon_F$；
\item 求基态总能量 $U$ 与二维“压强” $p$（即单位长度边界上的热力学压力对应的面密度表达）。
\end{enumerate}
\begin{solution}

二维动量空间中，每个量子态占据面积
\[
(2\pi/L)^2=\frac{(2\pi)^2}{A},
\]
自旋简并度为 $2$。半径不超过 $k$ 的圆盘中态数为
\[
N(k)=2\cdot \frac{A}{(2\pi)^2}\cdot \pi k^2=\frac{Ak^2}{2\pi}.
\]
由
\[
\varepsilon=\frac{\hbar^2k^2}{2m}
\]
得
\[
k^2=\frac{2m\varepsilon}{\hbar^2},
\]
故
\[
N(\varepsilon)=\frac{Am\varepsilon}{\pi\hbar^2}.
\]
因此态密度是常数：
\ansmath{
g(\varepsilon)=\frac{dN}{d\varepsilon}=\frac{Am}{\pi\hbar^2}.
}

令面密度
\[
n_s=\frac{N}{A},
\]
则在 $T=0$ 时，
\[
N=\int_0^{\varepsilon_F}g(\varepsilon)\,d\varepsilon
=\frac{Am}{\pi\hbar^2}\varepsilon_F.
\]
从而
\ansmath{
\varepsilon_F=\frac{\pi\hbar^2}{m}n_s.
}

基态总能量
\[
U=\int_0^{\varepsilon_F}\varepsilon\,g(\varepsilon)\,d\varepsilon
=\frac{Am}{\pi\hbar^2}\cdot \frac{\varepsilon_F^2}{2}.
\]
利用上式消去 $\varepsilon_F$，得到
\ansmath{
U=\frac12 N\varepsilon_F.
}

二维自由费米气体满足
\[
pA=U,
\]
故
\[
p=\frac{U}{A}=\frac12 n_s\varepsilon_F.
\]
代入 $\varepsilon_F$ 后
\ansmath{
p=\frac{\pi\hbar^2}{2m}n_s^2.
}
\end{solution}
\end{question}

